% =====================================================================
% AI Academy -- National AI Center
% Software Engineering Final Project -- Team Report Template
%
% HOW TO USE
% ----------
%   1. Edit the CONFIGURATION block below (team name, topic, members).
%   2. Replace every [TODO: ...] with your content. TODOs render in red,
%      so anything you forgot will be visually obvious in the PDF.
%   3. Three kinds of cues throughout:
%        - REQUIRED DELIVERABLE box (blue): exactly what we expect.
%        - PROBING QUESTIONS box (amber): NOT a checklist; these are
%          prompts that should shape how you write the section. A strong
%          response engages with at least 2 of them; a weak response
%          ignores them all.
%        - "Your response:" placeholders: replace with your text/diagram.
%   4. Target length: 8-10 pages including figures, excluding bibliography.
%   5. Compile with pdfLaTeX. Overleaf works out of the box.
% =====================================================================

\documentclass[11pt,a4paper]{article}

% ===== EARLY MACROS (must come before configuration) =================
% (Full styling lower down; we just need \TODO and \ph available
% before the configuration block uses them.)
\usepackage[dvipsnames]{xcolor}
\definecolor{accent2}{HTML}{8B0000}
\definecolor{ph}{HTML}{6B7280}
\newcommand{\TODO}[1]{\textcolor{accent2}{\textbf{[TODO: #1]}}}
\newcommand{\phh}[1]{\textcolor{ph}{\textit{#1}}}

% ===== CONFIGURATION (edit these) ====================================
\newcommand{\teamname}{Team \#1}
\newcommand{\topicchoice}{Topic 2 -- AI Food Analyzer}
\newcommand{\member}[2]{\textbf{#1} (#2)}  % name, GitHub handle
\newcommand{\reportdate}{19 September 2026}
\newcommand{\repolink}{\url{https://github.com/Gunas12/Food_Analyzer}}

% ===== course meta (rarely changes) ==================================
\newcommand{\coursecode}{AI-ENG-110}
\newcommand{\coursename}{Software Engineering}
\newcommand{\institution}{AI Academy, National AI Center}
\newcommand{\semester}{Spring 2026}

% ===== PACKAGES ======================================================
\usepackage[margin=2.0cm,headheight=15pt]{geometry}
\usepackage[T1]{fontenc}
\usepackage[utf8]{inputenc}
\usepackage[final,protrusion=true,expansion=false]{microtype}
\usepackage{amsmath,amssymb}
\usepackage{graphicx}
\usepackage{booktabs}
\usepackage{array}
\usepackage{tabularx}
\usepackage{ragged2e}
\usepackage{enumitem}
\usepackage[parfill]{parskip}
\setlength{\parskip}{0.55em}
\usepackage{listings}
\usepackage{tikz}
\usetikzlibrary{arrows.meta,positioning,fit,backgrounds}
\usepackage[most]{tcolorbox}
\usepackage{titlesec}
\usepackage{fancyhdr}
\usepackage{lastpage}
\usepackage[hidelinks]{hyperref}

\emergencystretch=3em
\hbadness=10000
\hfuzz=2pt

\hypersetup{
  colorlinks=true,
  linkcolor=NavyBlue,
  urlcolor=NavyBlue,
  citecolor=NavyBlue,
  pdftitle={\coursecode\ Final Project Report},
  pdfauthor={\teamname}
}

% ===== STYLING =======================================================
\definecolor{accent}{HTML}{0B3D91}
\definecolor{soft}{HTML}{F2F4F8}
\definecolor{deliver}{HTML}{E8F0FE}
\definecolor{deliverframe}{HTML}{1A56DB}
\definecolor{probe}{HTML}{FFF8E1}
\definecolor{probeframe}{HTML}{B45309}
\definecolor{ok}{HTML}{166534}

\titleformat{\section}
  {\Large\bfseries\color{accent}}{\thesection.}{0.6em}{}
\titleformat{\subsection}
  {\large\bfseries\color{accent}}{\thesubsection}{0.5em}{}

\newcommand{\ans}[1]{\rule[-2pt]{#1}{0.5pt}}
\let\ph\phh  % alias the early \phh as \ph for consistency with the body

% ---------------------------------------------------------------------
% GUIDANCE BOXES ON/OFF
% The template's blue "Required" and amber "Probing questions" boxes are
% instructions to us, not content. With them the PDF is 15 pages; the
% brief asks for 8-10. \guidancefalse hides them (the text still answers
% every one of them). Switch to \guidancetrue to see them again.
% ---------------------------------------------------------------------
\newif\ifguidance
\guidancefalse

% Deliverable callout (light blue)
\newtcolorbox{deliverboxshown}{
  colback=deliver, colframe=deliverframe, boxrule=0.5pt,
  arc=2pt, left=8pt, right=8pt, top=4pt, bottom=4pt
}
% Probing-question callout (light amber)
\newtcolorbox{probeboxshown}{
  colback=probe, colframe=probeframe, boxrule=0.5pt,
  arc=2pt, left=8pt, right=8pt, top=4pt, bottom=4pt
}

\ifguidance
  \newenvironment{deliverbox}{\begin{deliverboxshown}}{\end{deliverboxshown}}
  \newenvironment{probebox}{\begin{probeboxshown}}{\end{probeboxshown}}
\else
  \newenvironment{deliverbox}{\setbox0\vbox\bgroup}{\egroup\ignorespacesafterend}
  \newenvironment{probebox}{\setbox0\vbox\bgroup}{\egroup\ignorespacesafterend}
\fi

\newcommand{\answerbox}[1]{%
  \par\noindent
  \fcolorbox{black!30}{soft}{%
    \parbox{\dimexpr\linewidth-2\fboxsep-2\fboxrule\relax}{%
      \vspace{0pt}\centering
      \ph{[Your response here -- replace this placeholder.]}
      \vspace{#1}
    }%
  }\par
}

\newcommand{\figureplaceholder}[1]{%
  \par\noindent
  \begin{center}
  \fcolorbox{black!30}{soft}{%
    \parbox{0.75\linewidth}{%
      \vspace{0pt}\centering
      \ph{[Replace this box with \texttt{\textbackslash includegraphics\{...\}} or a tikzpicture.]}
      \vspace{#1}
    }%
  }
  \end{center}
}

\definecolor{codegreen}{rgb}{0,0.55,0}
\definecolor{codegray}{rgb}{0.4,0.4,0.4}
\definecolor{codepurple}{rgb}{0.55,0,0.78}
\definecolor{backcolour}{rgb}{0.97,0.97,0.97}

\lstdefinestyle{python}{
  language=Python,
  backgroundcolor=\color{backcolour},
  commentstyle=\color{codegreen},
  keywordstyle=\color{accent},
  numberstyle=\tiny\color{codegray},
  stringstyle=\color{codepurple},
  basicstyle=\ttfamily\footnotesize,
  breakatwhitespace=false,
  breaklines=true,
  keepspaces=true,
  numbers=left,
  numbersep=6pt,
  tabsize=2,
  frame=single,
  framerule=0.4pt,
  rulecolor=\color{black!30},
}

\pagestyle{fancy}
\fancyhf{}
\fancyhead[L]{\small\textit{\coursecode\ -- Final Report -- \teamname}}
\fancyhead[R]{\small\textit{\semester}}
\fancyfoot[L]{\small\textit{\institution}}
\fancyfoot[R]{\small Page \thepage\ of \pageref{LastPage}}
\renewcommand{\headrulewidth}{0.4pt}
\renewcommand{\footrulewidth}{0.4pt}

% =====================================================================
\begin{document}

\begin{center}
{\small \textsc{\institution} \\[0.2em] \textsc{\coursecode\ -- \coursename\ -- \semester}}\\[1.2em]
{\Large\bfseries\color{accent} Final Project Report}\\[0.3em]
{\LARGE\bfseries \topicchoice}\\[0.6em]
\rule{0.85\textwidth}{0.6pt}
\end{center}

\vspace{0.6em}
\begin{tabularx}{\textwidth}{@{}lXlX@{}}
\textbf{Team:} & \teamname & \textbf{Date:} & \reportdate \\
\textbf{Repo:} & \repolink & \textbf{Tag:} & \texttt{v1.0-final} \\
\textbf{Members:} &
\begin{tabular}[t]{@{}l l@{}}
\member{Gunash Mammadova}{@Gunas12} &
\member{Punhan Murselov}{@Punhan0} \\[0.2em]
\member{Ali Muradov}{@Ali-Muradov} &
\member{Esmira Mammadli}{@EsmiraMammadli}
\end{tabular}
& & \\
\end{tabularx}
\vspace{0.4em}\hrule\vspace{0.6em}

% ===== EXECUTIVE SUMMARY =============================================
\section*{Executive Summary}
\addcontentsline{toc}{section}{Executive Summary}

\begin{deliverbox}
\textbf{Required.} Exactly one paragraph (5--7 sentences). What you built, the headline concurrency speedup you measured, your headline robustness story, and the single most surprising thing the team learned. \emph{This is the first thing the grader reads. If you cannot write it clearly, you do not yet understand your own project.}
\end{deliverbox}

\begin{probebox}
\textit{Your paragraph should answer:} What problem did the system solve? What is the one number that proves the concurrency design works? What is the one failure mode you handled well? What would you change if you started over?
\end{probebox}

\par\noindent
We built a Food Analyzer that turns one meal photo into a calorie and
macronutrient breakdown: a vision-language model identifies ingredients with
estimated portions, and the USDA FoodData Central API supplies the nutrition
facts for each one. Our headline concurrency result is a \textbf{9.9$\times$}
speedup on a 10-ingredient nutrition lookup (1.517\,s sequential vs.\
0.154\,s concurrent, semaphore bound of 10), reproducible with
\texttt{python scripts/bench.py}. Our headline robustness story is that a
single ingredient's nutrition lookup can fail -- after its own retries are
exhausted -- without failing the whole analysis: \texttt{NutritionPipeline}
logs the miss and returns the meal with that one ingredient's macros at zero,
rather than discarding the entire request. The most surprising thing the
team learned while building this was how much of that robustness reduced to
a small number of narrowly-scoped \texttt{except ProviderError} clauses at
just two seams (the AI wrapper and the pipeline) rather than defensive
checks scattered throughout the codebase -- centralizing retries in one
\texttt{AIService} class turned out to be the single decision that made the
rest of the error handling simple. If we started over we would add a
per-call timeout (\texttt{asyncio.wait\_for}) inside \texttt{AIService}
from the first commit rather than last: our retry policy bounds how many
times a call can fail, but nothing bounds how long a single call may hang,
so the one failure mode we cannot currently survive is a provider that
accepts the connection and then never answers.

% ===== TABLE OF CONTENTS =============================================
\tableofcontents
\vspace{0.6em}\hrule\vspace{0.4em}

% =====================================================================
\section{Project Overview}

\subsection{Problem framing \& scope}

\begin{deliverbox}
\textbf{Required.} 2--3 paragraphs. State the topic problem in your own words. List the inputs, outputs, and any user-facing entry points (CLI commands, HTTP endpoints). State which provider stack you used (LLM, embedding, USDA, web search).
\end{deliverbox}

\begin{probebox}
\textit{Beyond the obvious:} which parts of the TOPIC.md spec did your team explicitly decide \emph{not} to do, and why? Where did you tighten the spec? Where did you loosen it?
\end{probebox}

\par\noindent
The system takes one meal photo as input and returns a structured breakdown
of its ingredients (name, estimated grams, confidence) together with
calories and macronutrients (protein, carbohydrates, fat) per ingredient and
as a meal total. It is exposed through two entry points that share the same
pipeline: a CLI (\texttt{python -m src.cli analyze <path>}) for local,
offline-friendly use, and an HTTP API (\texttt{POST /analyze},
\texttt{GET~/analyses/\{id\}}, \texttt{GET~/health}) for the browser demo
and for scripted access via \texttt{curl}. The provider stack is
configurable via \texttt{LLM\_PROVIDER} (\texttt{gemini} by default, with
\texttt{anthropic} and \texttt{openai} as drop-in alternatives, all through
the provided \texttt{ai/} package) for vision-language ingredient
identification, and the USDA FoodData Central API (\texttt{ai/nutrition.py})
for per-ingredient nutrition facts.

\par\noindent
We tightened the spec in two places the brief left open: image validation
checks magic bytes, not just the file extension or the client-supplied
\texttt{Content-Type}, so a mislabeled or corrupted file is rejected before
it ever reaches the VLM; and every uploaded image is written to disk under a
server-generated UUID filename rather than a client-supplied one, closing an
unrequested but obvious path-traversal risk. We loosened the spec, or left
it partially done, in a few places we consider defensible trade-offs under a
tight deadline rather than oversights: there is no \texttt{GET /analyses}
list endpoint, even though the underlying \texttt{Repository.list\_recent()}
method exists and is tested; there is no per-call timeout wrapping
\texttt{AIService} calls, only the retry/backoff policy (Section
\ref{sec:robustness} discusses this explicitly as a known gap rather than an
oversight discovered at grading time); and there is no failover for the
nutrition provider, only for the VLM.

\par\noindent
Two scope decisions came out of our own process rather than the brief. We
deliberately spent the last day of a three-day sprint closing gaps we had
found ourselves instead of starting an Advanced Bonus: \texttt{LOG\_LEVEL}
was declared in \texttt{Settings} but never applied to the root logger, and
the HTTP upload path trusted the client-supplied \texttt{Content-Type}
while the CLI path already checked magic bytes. Both were fixed (PR \#25,
\texttt{configure\_logging()} in \texttt{config.py} and the shared
\texttt{has\_image\_magic\_bytes()} helper in \texttt{validation.py}), and
we judged that closing a real inconsistency between our two entry points
was worth more than adding a feature neither entry point would have shared.
The second decision was to keep a single database table: a normalized
\texttt{ingredients} child table was discussed and rejected (Section
\ref{sec:storage}), because nothing in this topic ever queries one
ingredient across analyses.

\subsection{Provider choice and the AI module contract}

\begin{deliverbox}
\textbf{Required.} Name the LLM, embedding, and any auxiliary provider you used, and the model identifiers (\texttt{claude-sonnet-4-6}, \texttt{text-embedding-3-small}, etc.). State what \texttt{ai/} functions you call. \emph{Confirm in one sentence that you did not modify the public interface of the provided \texttt{ai/} package.}
\end{deliverbox}

\begin{probebox}
\textit{Push deeper:} how does your code defend against a malformed model response? What happens when the provider returns valid JSON that is semantically wrong (wrong topic enum, negative grams, etc.)?
\end{probebox}

\par\noindent
By default the system uses \texttt{LLM\_PROVIDER=gemini} with
\texttt{LLM\_MODEL=gemini-2.0-flash} for ingredient identification, with
\texttt{anthropic} and \texttt{openai} available as configuration-only
swaps through the same \texttt{ai.identify\_ingredients()} entry point.
Nutrition facts come from \texttt{NUTRITION\_PROVIDER=usda}, via
\texttt{ai.get\_nutrition\_provider()} and \texttt{USDAProvider.lookup()}.
We call exactly three functions from the provided package:
\texttt{ai.identify\_ingredients()}, \texttt{ai.get\_nutrition\_provider()}
/ \texttt{NutritionProvider.lookup()}, and \texttt{ai.compute\_totals()}. We
confirm that the public interface of \texttt{ai/} was never modified --
\texttt{git log --oneline -- ai/} shows the package touched only in the
commit that introduced it.

\par\noindent
The provided package already validates the model response's \emph{shape}:
\texttt{ai/schemas.py}'s Pydantic models reject a negative
\texttt{estimated\_grams}, an out-of-range \texttt{confidence}, or a missing
required field before the response ever reaches our code (the provided
\texttt{tests/test\_ai\_smoke.py} exercises exactly this, and we do not
weaken it). Our own layer adds two further defenses at the boundary: any
provider-side failure (malformed JSON, a 5xx, a timeout inside the USDA
client) is raised as a single \texttt{ProviderError} type that
\texttt{AIService} retries up to three times with exponential backoff before
giving up, and our own \texttt{IngredientLine}/\texttt{MealTotals} models
(\texttt{extra="forbid"}) reject any response shape that doesn't match
exactly what we expect to persist, rather than silently accepting and
storing unexpected extra fields.

% =====================================================================
\section{Architecture}\label{sec:arch}

\subsection{Module map}

\begin{deliverbox}
\textbf{Required.} An architecture diagram (boxes-and-arrows or sequence diagram) showing the main modules and their dependencies. Identify on the diagram: the boundary of the provided \texttt{ai/} package, your service layer that wraps it, your core business logic, your concurrency layer, your storage layer, and your CLI/HTTP entry points.
\end{deliverbox}

\begin{probebox}
\textit{Interpret your own diagram:} which arrows cross the AI-module boundary? Which arrows cross the storage boundary? If you had to swap providers (Anthropic $\to$ OpenAI), how many files would change?
\end{probebox}

\begin{center}
\begin{tikzpicture}[
  font=\footnotesize, scale=0.98, every node/.style={transform shape},
  box/.style={draw,rounded corners=2pt,minimum width=3.3cm,minimum height=0.8cm,align=center},
  entry/.style={box,fill=green!8,draw=ok,thick},
  sebox/.style={box,fill=accent!8,draw=accent},
  aibox/.style={box,fill=soft,draw=accent,very thick,minimum width=8.2cm},
  db/.style={box,fill=black!5,draw=black!50},
  arrow/.style={-Stealth,semithick},
  cross/.style={-Stealth,very thick,draw=accent}
]
  \node[entry] (cli) at (0,0)   {CLI --- \texttt{cli.py}\\\scriptsize\texttt{python -m foodanalyzer analyze}};
  \node[entry] (api) at (5.4,0) {HTTP API --- \texttt{api.py}\\\scriptsize\texttt{POST /analyze}};

  \node[sebox] (val)   at (-4.6,-1.6) {\texttt{validation.py}\\\scriptsize extension, size, magic bytes};
  \node[sebox] (core)  at (0,-1.6)    {\texttt{core/analyzer.py}\\\scriptsize FoodAnalyzer (CLI path)};
  \node[sebox] (lines) at (5.4,-1.6)  {\texttt{core/lines.py}\\\scriptsize shared response shaping};

  \node[sebox] (repo) at (-4.6,-3.2) {\texttt{storage/repository.py}};
  \node[sebox] (svc)  at (0.65,-3.2)    {\texttt{services/ai\_service.py}\\\scriptsize 3 attempts, backoff};
  \node[sebox] (pipe) at (5.4,-3.2)  {\texttt{concurrency/pipeline.py}\\\scriptsize gather + Semaphore(10)};

  \node[db]    (pg)    at (-4.6,-4.8) {PostgreSQL\\\scriptsize\texttt{analysis\_records}};
  \node[sebox] (cache) at (5.4,-4.8)  {\texttt{nutrition\_cache.py}\\\scriptsize TTL 24h, lock};

  \node[aibox] (ai) at (0.4,-6.4) {\textbf{provided \texttt{ai/} --- never modified}\\\scriptsize identify\_ingredients \ \textbullet\ \ NutritionProvider.lookup \ \textbullet\ \ compute\_totals};

  \draw[arrow] (cli) -- (core);
  \draw[arrow] (api) -- (lines);
  \draw[arrow] (core) -- (lines);
  \draw[arrow] (core) -- (val);
  \draw[arrow] (5.4,-0.4) -- (5.4,-0.85) -- (-4.6,-0.85) -- (val.north);
  \draw[arrow] (core) -- (svc);
  \draw[arrow] (lines) -- (pipe);
  \draw[arrow] (pipe) -- (svc);
  \draw[arrow] (pipe) -- (cache);
  \draw[cross] (svc) -- node[right,font=\scriptsize] {AI boundary} (ai);
  \draw[cross] (lines.east) -- (7.6,-1.6) -- (7.6,-6.4) -- node[above,pos=0.55,font=\scriptsize] {totals} (ai.east);
  \draw[arrow] (-1.3,-2.0) -- (-1.3,-3.2) -- (repo.east);
  \node[font=\scriptsize] at (-2.1,-2.65) {save / get};
  \draw[arrow] (7.3,0.4) -- (7.3,0.95) -- (-7.6,0.95) -- (-7.6,-3.2) -- (repo.west);
  \draw[arrow] (repo) -- (pg);
\end{tikzpicture}
\end{center}

\begin{center}
\phh{Green = entry points; blue = our SE layer (\texttt{src/}); heavy blue box = the
provided \texttt{ai/} package. Only the heavy arrows cross the AI boundary.}
\end{center}

\textbf{Module-by-module summary (1--2 sentences each):}

\begin{itemize}[leftmargin=1.4em,itemsep=1pt]
  \item \texttt{config.py} -- typed \texttt{Settings} (pydantic-settings), one cached \texttt{get\_settings()} instance for the whole process; also owns \texttt{configure\_logging()}, called once from each entry point to apply \texttt{LOG\_LEVEL} to the root logger.
  \item \texttt{validation.py} -- CLI-side image validation (existence, extension, size, magic bytes), plus \texttt{has\_image\_magic\_bytes()}, which \texttt{api.py} imports so the upload path rejects a spoofed \texttt{Content-Type} by exactly the same rule as the CLI rather than by a second, drifting copy of it.
  \item \texttt{wiring.py} -- composition root; builds \texttt{Repository} + \texttt{AIService} + \texttt{NutritionPipeline} from \texttt{Settings} for the CLI (the API builds an equivalent graph itself inside \texttt{create\_app()}).
  \item \texttt{core/analyzer.py} -- \texttt{FoodAnalyzer}, the one place the full pipeline (validate $\to$ identify $\to$ parallel nutrition lookup $\to$ totals $\to$ persist) is expressed as a single method for the CLI.
  \item \texttt{core/lines.py} -- \texttt{build\_ingredient\_lines()}, the shared seam between the CLI and the API: turns raw \texttt{Ingredient}/\texttt{NutritionFacts} into our own \texttt{IngredientLine}/\texttt{MealTotals}, so this specific piece of business logic (including the "missing lookup $\to$ zeroed macros, not a dropped ingredient" rule) cannot silently drift between the two entry points even though the surrounding orchestration is still written twice.
  \item \texttt{services/ai\_service.py} -- \texttt{AIService}, retries any \texttt{ai.*} call up to 3 times with exponential backoff on \texttt{ProviderError}.
  \item \texttt{services/nutrition\_cache.py} -- \texttt{NutritionCache}, a lock-protected, TTL-based in-memory cache keyed by normalized ingredient name.
  \item \texttt{concurrency/pipeline.py} -- \texttt{NutritionPipeline}, bounds concurrent nutrition lookups with \texttt{asyncio.Semaphore(MAX\_PARALLEL)} and consults the cache before every network call.
  \item \texttt{storage/repository.py} -- \texttt{Repository}, the only module that imports SQLAlchemy; everything else depends on its method signatures, not on the ORM.
  \item \texttt{api.py} / \texttt{cli.py} -- the two entry points; both validate, call the AI layer through \texttt{AIService}/
  \texttt{NutritionPipeline}, build the response through the shared \texttt{build\_ingredient\_lines()}, and persist -- but the five-step orchestration itself is still written once per entry point rather than both calling \texttt{FoodAnalyzer} (see Limitations).
\end{itemize}

\par\noindent
Only two arrows cross the \texttt{ai/} boundary -- \texttt{AIService} calling
\texttt{ai.identify\_ingredients()} and \texttt{NutritionPipeline} (via
\texttt{AIService}) calling a \texttt{NutritionProvider} -- which is exactly
why swapping \texttt{LLM\_PROVIDER} from \texttt{gemini} to
\texttt{anthropic} changes zero files in \texttt{src/}: it is a single
environment variable, since every provider implements the same
\texttt{VLMProvider} contract defined in the provided \texttt{ai/} package.
Only one arrow crosses the storage boundary -- \texttt{Repository.save()} /
\texttt{.get()} -- which is why both \texttt{tests/test\_repository.py} and
production code can point that one seam at either a live PostgreSQL
instance or an in-memory SQLite database without touching any other module.

\subsection{OOP design \& key abstractions}

\begin{deliverbox}
\textbf{Required.} Identify at least one example of inheritance or composition \emph{in your own code} (not the provided package's). Quote 5--10 lines of the relevant class or object design. Justify why this abstraction was the right tool; if you used an ABC, explain why.
\end{deliverbox}

\begin{probebox}
\textit{Defensibility:} if a grader asked "why is this an abstract base class rather than a Protocol, a plain function with a callable parameter, or a configuration enum?" --- what is your answer? Are there places where you used composition where inheritance would have been wrong, or vice versa?
\end{probebox}

\textbf{Code listing (composition, \texttt{src/core/analyzer.py}):}
\begin{lstlisting}[style=python]
class FoodAnalyzer:
    def __init__(
        self,
        ai_service: AIService,
        pipeline,  # NutritionPipeline
        repository: Repository,
        *,
        max_image_size_mb: float = 5.0,
    ) -> None:
        self._ai_service = ai_service
        self._pipeline = pipeline
        self._repository = repository
        self._max_image_size_mb = max_image_size_mb

    async def analyze(self, image_path: str) -> AnalysisRecord:
        ...  # validate -> identify -> parallel lookup -> totals -> save
\end{lstlisting}

\textbf{Justification:} \texttt{FoodAnalyzer} composes three collaborators
(\texttt{AIService}, \texttt{NutritionPipeline}, \texttt{Repository}) rather
than inheriting from a base pipeline class, because the four pipeline stages
are heterogeneous operations (an AI call, a fan-out over a network provider,
a database write) with no shared method signature an abstract base class
could usefully capture -- there is nothing to override, only things to
call in sequence. Composition also makes every collaborator independently
mockable in \texttt{tests/test\_analyzer.py} without a test double having to
satisfy an inheritance contract, which is the same reason
\texttt{src/storage/repository.py}'s \texttt{Repository} class wraps
SQLAlchemy behind three plain methods (\texttt{save}/\texttt{get}/\texttt{
list\_recent}) instead of subclassing a SQLAlchemy base: nothing else in the
codebase needs to be an ORM model, it only needs to call those three
methods. We did use inheritance once, indirectly, by depending on the
provided \texttt{ai/providers/base.py}'s \texttt{VLMProvider} ABC when
injecting a fake provider in tests (\texttt{ScriptedVLM(VLMProvider)}) --
there, unlike in our own code, an ABC is the right tool because every
provider genuinely shares one call signature (\texttt{describe(...)}) that
callers dispatch through polymorphically.

\subsection{Data flow across module boundaries}

\begin{deliverbox}
\textbf{Required.} List the dataclasses you defined for your own SE layer. Confirm in one sentence that no naked dictionaries cross module boundaries.
\end{deliverbox}

\begin{probebox}
\textit{The hardest call:} when did you reuse a dataclass from the provided \texttt{ai/} package vs.\ wrap it in a new one for your SE layer? What was the trigger for wrapping?
\end{probebox}

\par\noindent
Our SE layer defines four Pydantic models in \texttt{src/models.py}:
\texttt{IngredientLine} (one ingredient enriched with computed nutrition),
\texttt{MealTotals} (summed macros), \texttt{AnalysisResult} (what a fresh
analysis returns), and \texttt{AnalysisRecord} (the persisted row, with
optional \texttt{id}/\texttt{created\_at}). All four set
\texttt{model\_config = ConfigDict(extra="forbid")}, and every function
signature between \texttt{core/analyzer.py}, \texttt{api.py}, and
\texttt{storage/repository.py} is typed with these models or the provided
\texttt{ai.schemas} models -- no naked dictionary crosses a module boundary
anywhere in \texttt{src/}.

\par\noindent
The trigger for wrapping rather than reusing directly was persistence and
API-shape stability: we pass the provided \texttt{ai.schemas.Ingredient} and
\texttt{NutritionFacts} straight through \emph{inside} the pipeline (they
never touch the database), but the moment a value is about to leave our
layer -- returned from the HTTP API or written to PostgreSQL -- we convert
it into our own model. That boundary lets \texttt{IngredientLine} carry
fields the provided \texttt{Ingredient} does not have (computed
\texttt{kcal}/\texttt{protein\_g}/\texttt{carbs\_g}/\texttt{fat\_g} for this
ingredient's portion) without touching \texttt{ai/}, and lets our storage
schema evolve independently of whatever the provided package's schema does.

% =====================================================================
\section{Concurrency \& Performance}\label{sec:concurrency}

\subsection{Why this concurrency model}

\begin{deliverbox}
\textbf{Required.} State which concurrency primitive you used (\texttt{asyncio}, \texttt{ProcessPoolExecutor}, \texttt{ThreadPoolExecutor}) and why. Identify the bottleneck the parallel design targets. State the bound (e.g.\ \texttt{Semaphore(10)}) and how you picked it.
\end{deliverbox}

\begin{probebox}
\textit{Justify against alternatives:} could you have used threads instead of asyncio? Why not? What is the smallest workload at which the concurrent design starts to win against sequential? What happens if you set the semaphore to 1 vs.\ 100?
\end{probebox}

\par\noindent
We use \texttt{asyncio.gather} bounded by \texttt{asyncio.Semaphore(
MAX\_PARALLEL)} (default 10, configurable via \texttt{.env}) in
\texttt{NutritionPipeline.lookup()}. The bottleneck this targets is
network I/O wait: every ingredient's nutrition lookup is dominated by round
trip time to the USDA API (or, in the VLM step, to the configured LLM
provider), not by CPU work, so the event loop can issue the next lookup the
instant the current one starts waiting rather than blocking a whole OS
thread on it. \texttt{AIService} additionally runs every blocking provider
call through \texttt{asyncio.to\_thread}, so even a provider SDK with no
async support of its own does not block the event loop while it's inside a
retry. We picked 10 as the default bound because it matches the course
brief's sample workload sizes (10--20 ingredients/items) closely enough that
one batch clears in roughly one round trip, without picking a number large
enough to look like we were ignoring provider rate limits.

\par\noindent
Threads (a \texttt{ThreadPoolExecutor}) would work about as well for this
specific workload, since the work is I/O-bound rather than CPU-bound and
the GIL is released during network waits either way; we chose
\texttt{asyncio} because FastAPI's request handlers are already coroutines,
so the nutrition pipeline composes into that world (\texttt{await
pipeline.lookup(...)}) without a second concurrency model, and
\texttt{ProcessPoolExecutor} was never in consideration since there is no
CPU-bound work anywhere in this pipeline. The concurrent design starts
winning as soon as there are 2 ingredients to look up (any $N > 1$ with a
non-trivial round trip), since even two sequential HTTP calls take
noticeably longer than two that overlap. At \texttt{semaphore=1} the
pipeline is effectively sequential again -- lookups are still issued through
\texttt{asyncio.gather}, but the semaphore serializes them one at a time, so
the wall-clock time converges to the sequential baseline. At
\texttt{semaphore=100} with our default 10-ingredient benchmark workload the
result would not change at all, because the semaphore can never bind tighter
than the number of distinct ingredients in a batch -- the two numbers only
diverge on a meal with more than 100 ingredients, which is not a realistic
input for this topic.

\subsection{Sequential-vs-concurrent benchmark}

\begin{deliverbox}
\textbf{Required.} A table reporting wall-clock time for the same workload run sequentially vs.\ concurrently. State $N$ (number of requests), the machine, the python version, whether caches were cleared, and the command line. Reproduce the table also in your README.
\end{deliverbox}

\begin{probebox}
\textit{Honest interpretation:} did you reach the theoretical $N\times$ speedup? If not, what was the new bottleneck (provider rate limit? GIL? network RTT? single-threaded DB writes?). What would push the curve further?
\end{probebox}

\begin{center}\small
\begin{tabular}{lccccc}
\toprule
\textbf{Run} & $N$ & \textbf{Worker threads} & \textbf{Sequential (s)} & \textbf{Concurrent (s)} & \textbf{Speedup} \\
\midrule
(a) Windows dev machine (8 cores)      & 10 & 12 (default) & 1.517 & 0.154 & 9.9$\times$ \\
(b) 1-vCPU Linux container             & 10 & \phantom{0}5 (default) & 1.505 & 0.301 & 5.0$\times$ \\
(c) same container, executor raised    & 10 & 16 (forced)  & 1.505 & 0.152 & 9.9$\times$ \\
\bottomrule
\end{tabular}
\end{center}

\textit{Machine/version:} Windows, Python 3.13.15, run via
\texttt{\$env:PYTHONPATH="."; python scripts/bench.py}. \texttt{scripts/
bench.py} uses a fixed, simulated 150\,ms latency per lookup
(\texttt{\_SlowFakeProvider}) rather than a live USDA call, so the result
isolates the concurrency gain from \texttt{asyncio.gather} +
\texttt{Semaphore(10)} from provider-side variance; there is no cache to
clear since each run constructs a fresh \texttt{NutritionCache}.

\textbf{Discussion.} At $N=10$ and \texttt{max\_parallel=10} every lookup
fits in one batch, so the ceiling is close to $N\times$; run (a) reached
9.9$\times$, the missing 1\% being event-loop scheduling, not any provider
limit, since the benchmark never leaves the process.

\par\noindent
Re-running the identical command on a second machine is what made the
result interesting. On a 1-vCPU container the same code gave only
5.0$\times$ (run b), and the cause is not the semaphore. \texttt{AIService}
executes each blocking provider call through \texttt{asyncio.to\_thread},
which dispatches to the event loop's default \texttt{ThreadPoolExecutor},
sized by CPython at \texttt{min(32, os.cpu\_count() + 4)}. With one core
that is 5 threads, so five lookups run and the other five queue behind
them: $1.5\,\text{s}/0.3\,\text{s} = 5\times$, exactly the thread count.
Raising the executor to 16 workers on the same container restored
9.9$\times$ (run c), which confirms the mechanism rather than merely
suggesting it. The practical consequence is that
\texttt{Semaphore(MAX\_PARALLEL)} is only the \emph{intended} bound:
whenever \texttt{MAX\_PARALLEL} exceeds \texttt{cpu\_count + 4}, the real
bound is the thread pool, and on a small cloud instance (Render's free
tier included) a configured value of 10 quietly becomes 5. The fix is one
line at startup --
\texttt{loop.set\_default\_executor(ThreadPoolExecutor(max\_workers=settings.max\_parallel))}
-- and we would rather report the finding than a number that does not
transfer between machines.

\par\noindent
We did not re-run the benchmark against live USDA/Gemini backends, because
a live run would measure the providers' own variance and rate limits
rather than our concurrency design; we expect network RTT and USDA's
per-key limit to become the binding constraints there, but we have not
measured that and do not claim it.

\subsection{Rate limit \& politeness}

\begin{deliverbox}
\textbf{Required.} State the rate-limit strategy (sleep-on-429, token bucket, leaky bucket, request budget). Document the configured budget. State the highest burst your test suite produces.
\end{deliverbox}

\begin{probebox}
\textit{Failure-mode test:} have you actually observed a 429 from the provider during development, or is this strategy untested in anger? What does the log look like when the budget is hit?
\end{probebox}

\par\noindent
We do not implement a dedicated token-bucket or sleep-on-429 strategy;
politeness toward the provider comes from two mechanisms working together:
the \texttt{Semaphore(MAX\_PARALLEL)} bound (default 10) caps how many
nutrition lookups are in flight at once regardless of batch size, and
\texttt{AIService}'s generic retry policy backs off on any transient
failure without distinguishing rate limiting from other transient errors:
3 attempts, and therefore two waits, 0.5\,s then 1\,s, on any
\texttt{ProviderError} -- which is what both a 429 and a 5xx arrive as.
The configured budget is thus at most \texttt{MAX\_PARALLEL} (10) in-flight
nutrition calls and at most 3 provider calls per logical lookup. The
highest burst our test suite produces is bounded by the same semaphore:
\texttt{tests/test\_pipeline.py} parametrizes batches up to the configured
\texttt{max\_parallel}, and every provider call in tests is mocked, so no
test can trigger a live rate limit.

\par\noindent
To be explicit, because it matters for how much this section is worth: we
have \emph{never observed a real 429} from Gemini or USDA during
development. Our manual runs were single-image and well inside the free
tiers, so the rate-limit strategy above is reasoned, not battle-tested. We
did observe the neighbouring path -- a provider call failing repeatedly
and being retried, backed off, and mapped to a 503 (Section
\ref{sec:robustness}) -- and a 429 would travel exactly that path, but we
are not claiming to have seen the status code itself. The honest weakness
is that without parsing \texttt{Retry-After} we would back off for 1.5\,s
total where a provider might be asking for 60.

% =====================================================================
\section{Robustness \& Error Handling}\label{sec:robustness}

\subsection{Wrapping the AI module}

\begin{deliverbox}
\textbf{Required.} Describe your service-layer wrapper around \texttt{ai.*} calls: retries (which exceptions, how many attempts, backoff schedule), timeouts (per-call value), structured logging (what fields, at what level), input validation (what you check before calling, what you check on the response).
\end{deliverbox}

\begin{probebox}
\textit{Pressure-test:} what happens if the provider returns a 500? A 429? A connection reset mid-stream? A perfectly-formed JSON that contradicts the schema? Walk through one of these end-to-end.
\end{probebox}

\par\noindent
Every call to \texttt{ai.identify\_ingredients()} and every
\texttt{NutritionProvider.lookup()} goes through
\texttt{AIService.\_call\_with\_retry()}: it catches exactly
\texttt{ai.providers.base.ProviderError} (never a bare \texttt{except:}),
retries up to \texttt{max\_attempts=3} times with backoff starting at
\texttt{base\_delay=0.5}s and doubling (so two waits, 0.5s then 1s, before
the third and final attempt), and logs at \texttt{INFO} on every
start/retry/success -- explicitly \emph{without} the exception payload,
since a \texttt{ProviderError} can carry raw provider responses or partial
API keys (\texttt{tests/test\_ai\_service.py::
test\_logs\_record\_lifecycle\_without\_provider\_error\_payload} asserts
this). We do not currently set a per-call timeout in \texttt{AIService}
itself (Section~\ref{sec:limits}); the provided USDA client does set its
own internal 10s timeout. Input is checked before the call (magic bytes,
extension, size -- see the next subsection) and the response is checked
after: the provided \texttt{ai/schemas.py} rejects a malformed shape before
it reaches us, and \texttt{IngredientLine}/\texttt{MealTotals}
(\texttt{extra="forbid"}) reject anything we don't expect to persist.

\par\noindent
Walking through a 500/429 end-to-end: the provider SDK (inside
\texttt{ai/}) raises \texttt{ProviderError}; \texttt{AIService} catches it,
logs \texttt{"Retrying \{operation\} after attempt N/3 in D seconds"},
sleeps, and retries; if all 3 attempts fail, the original
\texttt{ProviderError} propagates to \texttt{api.py}, which maps it to an
HTTP \texttt{503} -- with a message that depends on \emph{which} call
exhausted its retries: \texttt{\{"detail": "AI provider unavailable."\}}
for a failed VLM identification, or \texttt{\{"detail": "Nutrition
provider unavailable."\}} for a failed USDA lookup, so the client (and
whoever is debugging) knows which half of the pipeline broke without
reading a stack trace. This is not a hypothetical: the exact first payload,
\texttt{artefacts/demo\_no\_meal.json}, is a live capture of this path from
our own manual testing --
\begin{lstlisting}[language=bash,numbers=none,basicstyle=\ttfamily\small]
$ curl -X POST http://localhost:8000/analyze -F "file=@data/no_meal_blue.png"
{"detail":"AI provider unavailable."}
\end{lstlisting}
-- no stack trace, no raw provider payload, reaches the caller. A
connection reset mid-stream inside the provider SDK is indistinguishable
from a 500 at our layer, since both surface as \texttt{ProviderError}; a
perfectly-formed JSON that contradicts the schema (e.g.\ negative grams) is
rejected one layer earlier, inside \texttt{ai/schemas.py}, before
\texttt{AIService} ever sees a result to retry.

\par\noindent
\textbf{The exact trigger}, since ``we observed X'' is worth more than
``the code should do X'': the capture comes from running the app against a
Gemini endpoint it could not reach --- the provider raises
\texttt{ProviderError} at the point where the client would otherwise talk
to \texttt{generativelanguage.googleapis.com}. The same state is
reproducible in one step by starting the API with no
\texttt{GOOGLE\_API\_KEY} in the environment or \texttt{.env}, since
\texttt{ai/providers/google.py} raises \texttt{ProviderError("GOOGLE\_API\_KEY
(or LLM\_API\_KEY) is not set.")} before any network call. Both routes
produce the same three-attempt sequence and the same 503; we re-ran it
while writing this section and the server log was:
\begin{lstlisting}[language=bash,numbers=none,basicstyle=\ttfamily\scriptsize]
INFO:src.services.ai_service:Starting identify_ingredients (attempt 1/3)
INFO:src.services.ai_service:Retrying identify_ingredients after attempt 1/3 in 0.500 seconds
INFO:src.services.ai_service:Starting identify_ingredients (attempt 2/3)
INFO:src.services.ai_service:Retrying identify_ingredients after attempt 2/3 in 1.000 seconds
INFO:src.services.ai_service:Starting identify_ingredients (attempt 3/3)
WARNING:src.api:vlm_call_failed: Gemini call failed: 403 Forbidden. {...}
INFO:httpx:HTTP Request: POST /analyze "HTTP/1.1 503 Service Unavailable"
\end{lstlisting}

\par\noindent
That last \texttt{WARNING} is worth stating precisely rather than
glossing, because it qualifies the claim above. \texttt{AIService} never
logs the exception payload; \texttt{api.py} logs it \emph{once}, at the
point where it converts the exhausted \texttt{ProviderError} into a 503,
because at that moment the operator has lost every other trace of what
went wrong. Nothing of it reaches the client: the HTTP body is exactly
\texttt{\{"detail":"AI provider unavailable."\}}. So the guarantee we
actually provide is ``no provider payload crosses the network boundary,
and none is written per-retry'' -- not ``no provider payload is ever
written''. If the deployment's logs were untrusted we would redact that
one line to the exception class name.

\subsection{Failure-mode analysis (one concrete incident)}

\begin{deliverbox}
\textbf{Required.} Pick \emph{one} failure mode you actually exercised (provider 5xx, timeout, malformed response, network drop, disk full, corrupted image, USDA outage, etc.). Describe: (i) how you injected the failure, (ii) what the system did, (iii) what the user saw, (iv) what the logs showed.
\end{deliverbox}

\begin{probebox}
\textit{Go beyond "graceful degradation worked":} what did you originally expect to happen vs.\ what actually happened? Did the test surface any design issue that you then fixed? Did the log give you enough to diagnose it without running the code again?
\end{probebox}

\par\noindent
\textbf{Injected failure:} a three-ingredient meal (\texttt{white rice},
\texttt{mystery sauce}, \texttt{grilled chicken}) run through
\texttt{NutritionPipeline.lookup()} against a scripted provider that always
raises \texttt{ProviderError} for \texttt{"mystery sauce"} (simulating a USDA
outage/no-match for one item) and succeeds normally for the other two.

\par\noindent
\textbf{What the system did:} \texttt{AIService} retried the failing lookup
3 times with backoff (0.05s, 0.1s in this run), then
\texttt{NutritionPipeline.\_lookup\_one} caught the exhausted
\texttt{ProviderError}, logged a \texttt{WARNING}, and returned
\texttt{None} for that one ingredient -- while \texttt{asyncio.gather} let
the other two lookups, which never touched the failing ingredient, complete
independently. The batch did not fail as a whole.

\par\noindent
\textbf{What the log showed} (captured directly, INFO level, unedited):
\begin{lstlisting}[language=bash,numbers=none,basicstyle=\ttfamily\scriptsize]
INFO:src.services.ai_service:Starting nutrition_lookup (attempt 1/3)
INFO:src.services.ai_service:Starting nutrition_lookup (attempt 2/3)
INFO:src.services.ai_service:Retrying nutrition_lookup after attempt 2/3 in 0.100 seconds
INFO:src.services.ai_service:Starting nutrition_lookup (attempt 3/3)
WARNING:src.concurrency.pipeline:Nutrition lookup failed for ingredient 'mystery sauce'
\end{lstlisting}

\par\noindent
\textbf{What the caller saw:} \texttt{white rice} $\to$ FOUND, \texttt{
mystery sauce} $\to$ MISSING (its \texttt{IngredientLine} carries
\texttt{kcal=0, protein\_g=0, carbs\_g=0, fat\_g=0} rather than being
dropped from the response), \texttt{grilled chicken} $\to$ FOUND -- the meal
total still reflects the two ingredients that did resolve. This matched
what we expected going in: the design goal was exactly that one unresolved
ingredient degrades that ingredient's numbers, not the whole request. The
log line alone (\texttt{"Nutrition lookup failed for ingredient
'mystery sauce'"}) is enough to know \emph{which} ingredient degraded
without re-running anything, though it does not say \emph{why} the
provider failed -- the underlying \texttt{ProviderError} message is
deliberately not logged (Section~\ref{sec:robustness} above), which is the
right trade-off for not leaking provider payloads but does mean a real
outage investigation would need to reproduce the failure with logging
temporarily turned up, or add a redacted error-class field to the warning.

\par\noindent
This scripted injection is the \emph{partial}-failure case (one ingredient
out of three). Its live counterpart, the \emph{total} failure of the VLM
half of the pipeline, is the capture described above: an unreachable
provider, three retries, and a 503 with no payload leak. We chose to
document the scripted one in detail because it is the more interesting
design decision -- deciding that a partial nutrition failure must still
return HTTP 200 with a degraded line is a choice, whereas returning 503
when the VLM is simply gone is not.

\subsection{Input validation}

\begin{deliverbox}
\textbf{Required.} For every entry point (CLI, HTTP, file, env), list the validation rules. State what error message the user gets when validation fails.
\end{deliverbox}

\begin{probebox}
\textit{Adversarial:} could a malformed image kill your service? A 100\,MB JSON request? A query of 50,000 characters? A path that escapes your storage directory?
\end{probebox}

\par\noindent
\textbf{CLI} (\texttt{validate\_image\_path}, checked in this order --
first failure wins): file must exist (\texttt{"file not found: \{path\}"});
extension must be \texttt{.jpg}/\texttt{.jpeg}/\texttt{.png}
(\texttt{"unsupported file extension ...; allowed: [...]"}\,); size must be
$\leq$ \texttt{MAX\_IMAGE\_SIZE\_MB} (\texttt{"file too large: X.X MB (max
Y MB)"}\,) and $>0$ bytes (\texttt{"file is empty"}); and the first 16 bytes
must match a real JPEG or PNG magic number (\texttt{"file content does not
look like a JPEG or PNG"}\,), so a text file renamed to \texttt{.png} is
still rejected. Any failure raises \texttt{ValidationError}, which
\texttt{cli.py} prints to \texttt{stderr} and exits with code \texttt{2}.

\par\noindent
\textbf{HTTP} (\texttt{\_validate\_and\_save\_image} in \texttt{api.py}):
\texttt{Content-Type} must be \texttt{image/jpeg} or \texttt{image/png}
(400); the upload must be non-empty (400) and $\leq$
\texttt{MAX\_IMAGE\_SIZE\_MB} (400); the on-disk filename is generated as a
fresh UUID rather than taken from the client (\texttt{\_safe\_upload\_path}
also asserts the resolved path still starts with the resolved upload
directory before writing), so a filename like \texttt{"../../etc/passwd"}
cannot write outside \texttt{uploads/} -- this defends the adversarial
case above directly rather than by accident.

\par\noindent
\textbf{Env:} \texttt{Settings} is a typed \texttt{pydantic-settings} model;
an out-of-range \texttt{postgres\_port}, a negative
\texttt{nutrition\_cache\_ttl\_seconds}, or an \texttt{LLM\_PROVIDER} outside
\texttt{\{anthropic, openai, gemini\}} fails at process startup with a
Pydantic validation error, not at request time.

\par\noindent
On the specific adversarial cases: a malformed image cannot kill the
service, since the magic-byte check rejects it before anything image-shaped
happens to it; a 100\,MB request is rejected as soon as its size is known
(400, before it is fully buffered against \texttt{MAX\_IMAGE\_SIZE\_MB}); we
have no free-text query parameter that accepts arbitrary length in this
topic (unlike Topics 3/4), so a 50{,}000-character string is not an
applicable attack surface here; and a path-escaping filename is defended
by the UUID-plus-resolved-path check described above.

\par\noindent
We did not want to leave that paragraph as reasoning about our own code,
so we ran each case against a live app instance (\texttt{TestClient} over
\texttt{create\_app()}, in-memory SQLite repository) and recorded the
actual responses:

\begin{center}\small
\begin{tabular}{p{6.6cm}cp{6.0cm}}
\toprule
\textbf{Request} & \textbf{Status} & \textbf{Body} \\
\midrule
100\,MB PNG upload & 400 & \texttt{File exceeds the 5 MB limit.} \\
\texttt{filename="../../etc/passwd.png"} & 503\footnotesize$^{*}$ & stored as \texttt{uploads/<uuid>.png}; no file created outside \texttt{uploads/} \\
Text bytes sent as \texttt{image/png} & 400 & \texttt{File content does not look like a JPEG or PNG.} \\
Zero-byte upload & 400 & \texttt{Uploaded file is empty.} \\
Valid JPEG sent as \texttt{text/plain} & 400 & \texttt{Unsupported content type: 'text/plain'.} \\
\bottomrule
\end{tabular}
\end{center}

{\footnotesize$^{*}$ The traversal attempt is not rejected -- it is
defused: the filename is discarded, the bytes are a valid PNG, so the
request proceeds normally and only then fails at the VLM call (the run
used an unreachable provider). The security-relevant assertion is the
second half of that cell: the file landed under \texttt{uploads/} with a
UUID name, and \texttt{../../etc/passwd.png} did not exist afterwards.}

% =====================================================================
\section{Storage \& Caching}\label{sec:storage}

\subsection{Schema}

\begin{deliverbox}
\textbf{Required.} The SQLite schema for your domain tables. Include indices and foreign keys. State which fields are derived (cached) vs.\ source-of-truth.
\end{deliverbox}

\begin{probebox}
\textit{Schema choices:} why did you store embedding bytes as a blob rather than serialized JSON, or vice versa? What happens to the DB if the embedding dimension changes (a different provider, a different model)?
\end{probebox}

We use PostgreSQL (not SQLite) as the source of truth, via a SQLAlchemy
async engine (\texttt{src/storage/repository.py}); tests run the identical
schema against \texttt{sqlite+aiosqlite:///:memory:} so the suite stays
offline. Equivalent DDL for the one table, \texttt{analysis\_records}:
\begin{lstlisting}[style=python,language=sql,numbers=none,basicstyle=\ttfamily\small]
CREATE TABLE analysis_records (
  id               SERIAL PRIMARY KEY,        -- autoincrement
  created_at       TIMESTAMPTZ NOT NULL,      -- default set in Python, not SQL
  image_path       TEXT NOT NULL,
  meal_recognized  BOOLEAN NOT NULL DEFAULT TRUE,
  ingredients      JSON NOT NULL,   -- list[IngredientLine], denormalized
  totals           JSON NOT NULL   -- MealTotals, derived from ingredients
);
\end{lstlisting}
There are no foreign keys or secondary indices -- this is a single-table
history log with one access pattern (\texttt{list\_recent()} orders by
\texttt{created\_at DESC, id DESC}), so we did not add an index beyond the
primary key given the project's scale. \texttt{ingredients} and
\texttt{totals} are derived/denormalized: both are recomputed from the VLM
response and USDA lookups at analysis time and stored as JSON for fast
reads, rather than normalized into a separate \texttt{ingredients} table --
the source of truth for any individual value is the analysis run itself,
not the row. One clarification on the defaults: \texttt{id} is a real
autoincrement column, but \texttt{created\_at} and the JSON columns get
their defaults from the SQLAlchemy model in Python rather than from SQL
\texttt{DEFAULT} clauses, so a row inserted by anything other than our
\texttt{Repository} would need to supply them.

\par\noindent
This topic has no embedding step, so the blob-versus-JSON question does not
apply directly; the equivalent decision we did make is \textbf{JSON columns
versus a normalized \texttt{ingredients} child table} with a foreign key
back to \texttt{analysis\_records}. We chose JSON because of the access
pattern: every read in the system fetches one complete analysis
(\texttt{get(id)}) or the last $N$ analyses (\texttt{list\_recent()}), and
nothing ever queries across ingredients -- there is no ``how many times has
this user eaten rice'' feature, and inventing the join table for a query
nobody makes would cost an extra round trip on every read plus insert
ordering we would have to get right. The cost of that choice is explicit:
the moment such a query is wanted, the data has to be migrated out of JSON,
and PostgreSQL cannot enforce the shape of what we store -- the
\texttt{extra="forbid"} Pydantic models are the only thing validating it,
which is precisely why they exist. A schema change on our side
(e.g.\ adding fibre to \texttt{IngredientLine}) also leaves old rows with
the old shape, so a reader must tolerate both; today it does, because rows
are only ever read back through the same model version that wrote them.

\subsection{Caching policy}

\begin{deliverbox}
\textbf{Required.} What do you cache, how long, where, and how is the cache invalidated. Report cache hit rate from a demo run (a single number is fine).
\end{deliverbox}

\begin{probebox}
\textit{Correctness vs.\ freshness:} which cached values can become stale (e.g.\ USDA nutrition for a renamed food, an LLM summary if the prompt changes)? How would you detect and evict them in production?
\end{probebox}

\par\noindent
We cache USDA nutrition lookups only (not VLM identifications, which depend
on the specific photo and are never repeated). \texttt{NutritionCache}
(\texttt{src/services/nutrition\_cache.py}) is an in-memory,
\texttt{threading.Lock}-protected dictionary keyed by
\texttt{normalize\_ingredient\_name(name)} (stripped, case-folded), with a
default TTL of \texttt{NUTRITION\_CACHE\_TTL\_SECONDS=86400} (24h) tracked
per-entry via a monotonic clock; an expired entry is evicted lazily, on the
next \texttt{get()} for that key, rather than by a background sweep. The
lock exists because \texttt{AIService} runs each blocking provider call
through \texttt{asyncio.to\_thread}, so concurrent nutrition lookups can
genuinely touch the cache from different OS threads, not just interleaved
coroutines. On a demo run of 5 meals sharing a common set of ingredients (15 ingredient
lookups over 8 distinct ingredients) we measured a \textbf{46.7\% cache hit
rate} -- 7 of 15 lookups served from cache, 8 real provider calls. It is
reproducible in one command, with no API keys, via
\texttt{scripts/cache\_demo.py}, which drives five meals through one shared
\texttt{NutritionPipeline} against a provider that counts its own calls:
\begin{lstlisting}[language=bash,numbers=none,basicstyle=\ttfamily\small]
$ PYTHONPATH=. python scripts/cache_demo.py
ingredient lookups          15
provider calls               8
cache hits                   7
cache hit rate: 46.7%
\end{lstlisting}
Two distinct mechanisms produce that number and we keep them separate: the
\texttt{NutritionCache} serves repeats \emph{across} meals, while
\texttt{NutritionPipeline.lookup()} additionally collapses duplicate
normalized names \emph{within} one batch before dispatching, so the same
ingredient listed twice in one photo is one provider call even with the
cache disabled.

\par\noindent
The cached value can go stale if USDA revises a food's nutrition data
within the 24h window (rare, but not impossible for a database that is
itself periodically updated), or if the VLM returns a different string for
the same real ingredient across two photos (\texttt{"grilled chicken"} vs.\
\texttt{"grilled chicken breast"}), which our exact-string-after-normalize
key would treat as two different cache entries rather than detecting they
are the same food. In production we would not change the TTL-based
eviction itself, but would add a manual or admin-triggered
\texttt{cache.clear()} keyed to a known USDA data-refresh schedule, and
consider fuzzy-matching or synonym normalization for the second problem
before it can be solved by caching policy alone.

% =====================================================================
\section{Testing}\label{sec:testing}

\subsection{Strategy \& coverage}

\begin{deliverbox}
\textbf{Required.} Report the coverage number (\texttt{pytest --cov}). State how you mocked the AI module and the HTTP layer. Confirm the provided \texttt{tests/test\_ai\_smoke.py} still passes.
\end{deliverbox}

\begin{probebox}
\textit{Quality vs.\ coverage:} which lines did you decide \emph{not} to cover, and why? Which mocks did you have to redesign mid-project because they were testing the wrong thing?
\end{probebox}

\begin{center}\small
\begin{tabular}{lcc}
\toprule
\textbf{Suite} & \textbf{Tests} & \textbf{Result} \\
\midrule
Our own tests, 11 files (\texttt{--cov=src}) & 114 & pass, \textbf{97\%} coverage \\
Provided \texttt{ai/} smoke tests (untouched) & \phantom{0}29 & pass \\
\midrule
\texttt{python -m pytest -q} (whole suite) & \textbf{143} & 143 passed \\
\bottomrule
\end{tabular}
\end{center}

\par\noindent
Per-module breakdown (\texttt{pytest --cov=src --cov-report=term-missing},
490 statements total, 17 missed, exact output archived in
\texttt{artefacts/coverage.txt}): every module in \texttt{src/} is at 100\%
except \texttt{api.py} (118 stmts, 15 missed, \textbf{87\%} -- lines 58, 65,
83, 92, 140, 142, 156, 169--171, 185--187, 205--206), \texttt{cli.py} (42
stmts, 1 missed, \textbf{98\%} -- line 82), and
\texttt{storage/repository.py} (44 stmts, 1 missed, \textbf{98\%} -- line
128). The AI module is mocked via a fake \texttt{VLMProvider} subclass
(\texttt{ScriptedVLM}, returning or raising scripted outcomes in order)
injected through \texttt{AIService}'s \texttt{vlm=} parameter, and a fake
\texttt{NutritionProvider} for nutrition lookups -- no test imports a real
provider SDK. The HTTP layer is exercised with FastAPI's
\texttt{TestClient} against \texttt{create\_app()} built with an in-memory
SQLite repository and fake AI/nutrition factories, so \texttt{tests/
test\_api.py} never opens a socket or hits PostgreSQL. \texttt{tests/
test\_ai\_smoke.py} (the 29 provided tests) is untouched and passes.

\par\noindent
The 17 uncovered statements are almost entirely defensive branches that are
straightforward to reason about but awkward to trigger without a real ASGI
server or a real terminal (e.g.\ the static-file-mount
\texttt{if static\_dir.is\_dir()} guard, and a couple of \texttt{except}
branches around provider construction at request
time) -- we judged the risk/effort trade-off of covering them as low given
the deadline.

\par\noindent
One mock did have to be redesigned mid-project. The first fake VLM returned
a fixed list of ingredients on every call, which was fine for the happy
path but useless for the thing we actually needed to prove: that attempt 2
succeeds after attempt 1 raises. A mock that cannot distinguish call 1 from
call 2 cannot test a retry policy at all -- it only tests that the wrapper
does not crash. We replaced it with \texttt{ScriptedVLM}, which is handed a
\emph{sequence} of outcomes (values or exceptions) and consumes one per
call, and the same shape was then reused for the fake nutrition provider,
which is what makes the partial-failure pipeline test in
Section~\ref{sec:robustness} expressible. A second, smaller correction:
\texttt{AIService} takes an injectable \texttt{sleep} callable, added once
we noticed the retry tests were spending real wall-clock seconds sleeping
through the production backoff schedule.

\subsection{Notable tests}

\begin{deliverbox}
\textbf{Required.} Highlight at least three tests: (i) one happy-path end-to-end test, (ii) one concurrency test (e.g.\ that \texttt{asyncio.gather} behaves correctly when one task raises), (iii) one error-path test that surfaces a real bug you would otherwise have missed.
\end{deliverbox}

\begin{probebox}
\textit{What didn't you test?} What is the test you wish you had time to write?
\end{probebox}

\begin{itemize}[leftmargin=1.4em,itemsep=3pt]
  \item \textbf{Happy path:} \texttt{tests/test\_api.py::
  test\_analyze\_happy\_path} -- POSTs a real sample image through
  \texttt{create\_app()}'s in-memory stack and asserts the full response
  shape (ingredients, totals, persisted \texttt{id}); mirrored at the
  business-logic level by \texttt{tests/test\_analyzer.py::
  test\_happy\_path\_recognized\_meal}.
  \item \textbf{Concurrency:} \texttt{tests/test\_pipeline.py::
  test\_provider\_failure\_preserves\_other\_results
  \_and\_logs\_warning}
  -- runs a batch where one ingredient's lookup raises, and asserts the
  other ingredients' results are still returned via \texttt{asyncio.gather}
  (this is the test behind the failure-mode walkthrough in
  Section~\ref{sec:robustness}); \texttt{test\_max\_parallel\_bounds\_calls\_
  across\_concurrent\_batches} additionally asserts the semaphore bound is
  never exceeded, across parametrized limits.
  \item \textbf{Error path (surfaced a real design question):}
  \texttt{tests/test\_api.py::test\_analyze\_unknown
\_ingredient\_still\_
  returns\_ok} -- forced us to decide, and then assert, that a nutrition
  lookup failure returns HTTP \texttt{200} with a degraded ingredient
  line rather than a \texttt{503} for the whole request; without this test
  it would have been easy to let a single missing ingredient fail the
  entire analysis instead.
\end{itemize}

\par\noindent
\textbf{What we did not test.} Three gaps, in the order we would close
them. First, no test exercises a hung provider, because there is no
timeout to exercise -- the test and the feature are the same missing piece
(Section~\ref{sec:limits}); the test we would write first is a fake
provider that blocks forever and an assertion that \texttt{/analyze}
answers within $n$ seconds, which today would simply hang the suite.
Second, nothing tests \emph{concurrent requests}: every concurrency test
we have lives inside one request, fanning out over ingredients, so the
interesting question -- what happens when ten clients each open a batch of
ten lookups against one shared cache and one shared \texttt{AIService} --
is unmeasured. Third, no test runs against a real PostgreSQL instance;
the suite uses SQLite, which means a PostgreSQL-specific problem (JSON
type handling, or a connection pool exhausted under load) would reach
production untested. With another day, the concurrent-request test is the
one we would write, because it is the scenario a grader is most likely to
create by accident during the demo.

% =====================================================================
\section{Deployment \& Reproducibility}\label{sec:deploy}

\subsection{Docker image}

\begin{deliverbox}
\textbf{Required.} The exact \texttt{docker build} and \texttt{docker run} commands. The image size. The Python version baked in. Note whether you use a single-stage or multi-stage build.
\end{deliverbox}

\begin{probebox}
\textit{Engineering trade-offs:} did you ship an Alpine, slim, or full base image, and why? How does your image size compare to a fresh \texttt{python:3.12-slim} ($\sim$120\,MB)?
\end{probebox}

\par\noindent
\begin{lstlisting}[language=bash,numbers=none]
$ docker compose up --build   # builds the API image and starts it with
PostgreSQL
$ docker build -t food-analyzer .
$ docker run -p 8000:8000 --env-file .env food-analyzer
\end{lstlisting}
Python \texttt{3.12-slim} is baked in, and the build is multi-stage: a
\texttt{builder} stage creates a virtualenv and installs
\texttt{requirements.txt} into it, and the runtime stage copies only
\texttt{/opt/venv} plus \texttt{ai/}, \texttt{data/}, \texttt{src/},
\texttt{static/} out of it -- \texttt{pip}'s build cache and any compiler
toolchain used to build wheels never reach the final image. The container
also runs as a non-root \texttt{appuser}. We chose \texttt{slim} over
\texttt{alpine} because several dependencies (\texttt{asyncpg},
\texttt{cryptography} via \texttt{google-auth}) ship prebuilt \texttt{
manylinux} wheels for glibc but not always for musl, so \texttt{alpine}
risks a slower or failing build for a marginal size win; we chose it over
the full \texttt{python:3.12} image because we don't need a system compiler
or dev headers in the runtime stage.

\par\noindent
\textbf{Size, and an honest look at it.} The image is roughly
\textbf{400\,MB}: the \texttt{python:3.12-slim} base ($\sim$130\,MB), the
copied \texttt{/opt/venv} (\textbf{273\,MB} measured by installing
\texttt{requirements.txt} into a clean virtualenv), and a few MB of
application code and sample images. So we are about 3$\times$ the bare
base image, and essentially all of that is dependencies rather than
anything we wrote. The largest single items are \texttt{numpy}
($\sim$70\,MB with its bundled libraries), \texttt{sqlalchemy}
($\sim$23\,MB), \texttt{cryptography} and \texttt{asyncpg}
($\sim$15\,MB each). The finding we did not expect is that
\textbf{$\sim$70--80\,MB of the runtime image is test and type-checking
tooling}: \texttt{requirements.txt} is a single list, so \texttt{mypy}
(including its compiled \texttt{mypyc} extension), \texttt{pytest} and
\texttt{pygments} are installed into the venv that the runtime stage
copies, even though the container never runs a test. Splitting
\texttt{requirements-dev.txt} out would cut close to a fifth of the image
for no behavioural change, and that is the first thing we would do to this
Dockerfile. A second, smaller defect from the same reading: the runtime
stage copies \texttt{ai/}, \texttt{data/}, \texttt{src/} and
\texttt{static/} but not the \texttt{foodanalyzer/} entrypoint package, so
\texttt{python -m foodanalyzer} works from a checkout but not inside the
container, where the CLI must be invoked as \texttt{python -m src.cli}.

\subsection{Live deployment}

\par\noindent
Beyond the local Docker stack, the API and web UI are deployed on Render
(free tier) at \url{https://food-analyzer-tx9l.onrender.com/}, with
interactive Swagger UI docs at \texttt{/docs} and a health check at
\texttt{/health} -- this is the same container image described above, not
a separate build. The free tier spins the service down after a period of
inactivity, so the first request after idle time can take 30--60s to wake
it back up; this is a Render platform characteristic, not an application
bug, but is worth mentioning during the oral defense demo in case the
grader hits a cold start.

\subsection{Configuration}

\begin{deliverbox}
\textbf{Required.} A short list of every environment variable the app reads, and what happens if each is missing.
\end{deliverbox}

\par\noindent
Every variable below has a default in \texttt{Settings} except the three
API keys, so a missing \texttt{.env} does not crash the process at startup
-- it fails the first time a request actually needs the missing value.

\begin{center}\small
\begin{tabular}{lll}
\toprule
\textbf{Variable} & \textbf{Default} & \textbf{If missing} \\
\midrule
\texttt{LLM\_PROVIDER} & \texttt{gemini} & uses default provider \\
\texttt{LLM\_MODEL} & \texttt{gemini-2.0-flash} & uses default model \\
\texttt{GOOGLE\_API\_KEY} / \texttt{ANTHROPIC\_API\_KEY} / & \texttt{None} & first VLM call raises \\
\quad\texttt{OPENAI\_API\_KEY} & & \texttt{ProviderError} \\
\texttt{USDA\_API\_KEY} & \texttt{None} & first nutrition call raises \\
\texttt{POSTGRES\_HOST/PORT/DB/USER/PASSWORD} & see README & connects to local defaults \\
\texttt{NUTRITION\_CACHE\_TTL\_SECONDS} & 86400 & 24h cache TTL \\
\texttt{MAX\_IMAGE\_SIZE\_MB} & 5 & 5\,MB upload limit \\
\texttt{MAX\_PARALLEL} & 10 & semaphore bound of 10 \\
\texttt{HTTP\_PORT} & 8000 & API listens on 8000 \\
\texttt{LOG\_LEVEL} & \texttt{INFO} & root logger stays at \texttt{INFO} \\
\texttt{UPLOAD\_DIR} & \texttt{uploads} & uploads written to \texttt{./uploads} \\
\bottomrule
\end{tabular}
\end{center}

% =====================================================================
\section{Limitations \& Future Work}\label{sec:limits}

\begin{deliverbox}
\textbf{Required.} 3--5 bullet points. Honest list of where the system is brittle, what it would take to productionize, and what you would do differently with another week.
\end{deliverbox}

\begin{probebox}
\textit{Be honest:} no project is perfect. What was the worst design decision you made? What did you not get to test under load? What is the scenario most likely to break the system on a busy day?
\end{probebox}

\begin{itemize}[leftmargin=1.4em,itemsep=3pt]
  \item \textbf{No per-call timeout in \texttt{AIService}.} Retries are
  bounded (3 attempts, exponential backoff), but nothing caps how long a
  single call to \texttt{ai.identify\_ingredients()} or a nutrition lookup
  can block -- a hung VLM call would hang the whole request indefinitely.
  Only the provided USDA client has its own internal 10s timeout.
  \item \textbf{\texttt{GET /analyses} (list) does not exist.}
  \texttt{Repository.list\_recent()} is implemented and tested, but no
  route calls it, so the API can create and fetch single records but not
  browse history.
  \item \textbf{The API and CLI still write the orchestration twice.}
  \texttt{cli.py} calls \texttt{FoodAnalyzer.analyze()}, while
  \texttt{api.py} builds an equivalent five-step flow inline inside
  \texttt{create\_app()} instead of reusing that class -- we did close half
  of this gap by pulling the response-shaping step
  (validated ingredient $\to$ \texttt{IngredientLine}/\texttt{MealTotals},
  including the "missing lookup degrades to zero, doesn't drop the
  ingredient" rule) out into the shared \texttt{core/lines.py}, so that one
  piece of business logic cannot drift between the two paths even though
  the surrounding five-step sequence is still written once per entry point.
  \item \textbf{Single nutrition provider, no failover.} If USDA is down,
  every analysis in that window returns ingredients with zeroed macros
  rather than falling back to a second provider -- there is no equivalent
  of the \texttt{LLM\_PROVIDER} swap-by-env-var for nutrition data.
\end{itemize}

\begin{itemize}[leftmargin=1.4em,itemsep=3pt]
  \item \textbf{\texttt{MAX\_PARALLEL} is not the real concurrency bound.}
  Because provider calls run through \texttt{asyncio.to\_thread}, the
  effective limit is the smaller of the semaphore and the default thread
  pool, \texttt{min(32, cpu\_count + 4)} -- 5 on a single-core host
  (Section~\ref{sec:concurrency}). The configuration therefore promises
  more than it delivers on exactly the kind of small instance we deployed
  to.
  \item \textbf{The cache is per-process and in-memory.} Two Uvicorn
  workers keep two caches with no shared state, and every restart is a
  cold start -- fine at our scale, wrong the moment the service is scaled
  horizontally, where Redis with the same TTL would be the replacement.
\end{itemize}

\par\noindent
\textbf{The worst decision, in hindsight}, was writing the API's
orchestration inline in \texttt{create\_app()} instead of calling
\texttt{FoodAnalyzer}. It felt cheaper on day 2 because the HTTP path
needed upload handling and per-step 503s that the CLI path did not, and
every later change then had to be made twice. We caught one instance of
the resulting drift (the zeroed-macros rule) and fixed it by extracting
\texttt{core/lines.py}, which closed roughly half the gap; we did not
finish the job. The lesson we would carry forward is that the right move
was to make \texttt{FoodAnalyzer} accept already-validated bytes and let
each entry point own only its own I/O, rather than to duplicate the
sequence and clean up afterwards.

\par\noindent
\textbf{What we never tested under load:} anything involving more than one
request at a time. Our concurrency work is all \emph{within} a request --
ten ingredients fanned out inside one \texttt{/analyze} call -- and we ran
no test, benchmark, or manual trial with several clients posting at once.
So the numbers in Section~\ref{sec:concurrency} describe a single user
with a busy plate, not a busy service.

\par\noindent
\textbf{The scenario most likely to break us on a busy day} follows
directly from that gap: $k$ simultaneous \texttt{POST /analyze} requests,
each opening up to 10 nutrition lookups, produce up to $10k$ in-flight
provider calls, because the semaphore is created per pipeline instance and
bounds one request, not the process. Ten concurrent users is a plausible
demo-day load and 100 outstanding USDA calls is not polite; the first
symptom would be USDA 429s, which our retry would absorb silently into
slower responses rather than reporting as rate limiting. The fix is small
and we know what it is -- one process-wide semaphore created at startup and
shared by every request, instead of one per pipeline -- but it is not in
the code we are submitting.

% =====================================================================
\section{Tools \& Acknowledgements}\label{sec:tools}

\begin{deliverbox}
\textbf{Required.} Disclose substantial AI-assistant use. For each significant LLM-aided contribution: which module, which assistant (Claude/ChatGPT/Cursor/Copilot/etc.), and whether the team reviewed and adapted it.
\end{deliverbox}

\begin{probebox}
\textit{Cultural note:} "AI helped with everything" is not a useful disclosure; "Cursor drafted the initial \texttt{retry\_policy.py}; we rewrote backoff logic after observing it hit the rate limit" is.
\end{probebox}

 
 

\par\noindent
Claude was used throughout the project as a development and documentation assistant. It helped draft the \texttt{README.md}, architecture and API documentation, the initial report draft, and parts of the implementation, particularly around retry logic, asynchronous concurrency, SQLAlchemy patterns, and the frontend interface.

All AI-generated material was reviewed, tested, and adapted by the team. We corrected factual issues, verified the implementation against the actual codebase and test results, and made the final design and coding decisions ourselves. All tests were written manually by the respective module owners. Each team member understands the complete codebase and can explain and defend the implementation during the oral defense.




% =====================================================================
\section*{References}
\addcontentsline{toc}{section}{References}

\begin{enumerate}[label={[\arabic*]},leftmargin=2.4em,itemsep=2pt]
  \item FastAPI documentation -- \url{https://fastapi.tiangolo.com/}
  \item Pydantic Settings documentation -- \url{https://docs.pydantic.dev/latest/concepts/pydantic_settings/}
  \item SQLAlchemy 2.0 Asyncio documentation -- \url{https://docs.sqlalchemy.org/en/20/orm/extensions/asyncio.html}
  \item Python \texttt{asyncio} documentation, \texttt{Semaphore} and \texttt{gather} -- \url{https://docs.python.org/3/library/asyncio-sync.html}
  \item USDA FoodData Central API documentation -- \url{https://fdc.nal.usda.gov/api-guide}
  \item M. Brooker, \emph{Timeouts, retries, and backoff with jitter} -- Amazon Builders' Library, \url{https://aws.amazon.com/builders-library/timeouts-retries-and-backoff-with-jitter/}. Shaped the retry policy: bounded attempts and exponential backoff, and the reason our lack of a per-call timeout is listed as the first limitation rather than a detail.
  \item Python documentation, \texttt{asyncio.to\_thread} and \texttt{loop.set\_default\_executor} -- \url{https://docs.python.org/3/library/asyncio-task.html\#asyncio.to_thread}. The default-executor sizing documented here is what explains the 5.0$\times$ vs.\ 9.9$\times$ benchmark result in Section~\ref{sec:concurrency}.
  \item OWASP, \emph{File Upload Cheat Sheet} -- \url{https://cheatsheetseries.owasp.org/cheatsheets/File_Upload_Cheat_Sheet.html}. Source of the two upload rules we added beyond the brief: never reuse the client-supplied filename, and validate content rather than \texttt{Content-Type}.
\end{enumerate}

\appendix
% =====================================================================
\section{Appendix --- Reproducing the benchmark}
\textbf{Exact commands to reproduce the §\ref{sec:concurrency} benchmark.} Anyone with the repo and a valid \texttt{.env} should be able to run the lines below and get numbers within 20\% of the ones in the table.

\begin{lstlisting}[style=python,language=bash,numbers=none]
$ git clone https://github.com/Gunas12/Food_Analyzer.git
$ cd Food_Analyzer
$ python -m venv .venv && .venv\Scripts\activate   # or source .venv/bin/activate
$ pip install -r requirements.txt
$ $env:PYTHONPATH="."                              # or export PYTHONPATH=.
$ python scripts/bench.py
\end{lstlisting}

\texttt{scripts/bench.py} uses a scripted \texttt{\_SlowFakeProvider} with a
fixed simulated latency, not a live USDA call, so this reproduces without
any API keys or \texttt{.env} -- the same command was used to produce the
9.9$\times$ figure in Section~\ref{sec:concurrency}.

\vspace{1em}\hrule\vspace{0.6em}
\noindent\small\textit{End of report --- thank you for reading.}

\end{document}